\documentclass[aspectratio=169,10pt]{beamer}
% Shared Beamer setup for Fall 2026 course slides.
% Clean Cal Poly-inspired theme: no custom class files or uncommon packages.
\mode<presentation>{
  \usetheme{default}
  \usefonttheme{professionalfonts}
}

\definecolor{CalPolyGreen}{HTML}{154734}
\definecolor{CalPolyGold}{HTML}{BD8B13}
\definecolor{CalPolySage}{HTML}{7FA28F}
\definecolor{CalPolyMint}{HTML}{DDF1E7}
\definecolor{CalPolySand}{HTML}{A29F86}
\definecolor{CalPolyBeige}{HTML}{EFEEDE}
\definecolor{DeckInk}{HTML}{1F2933}
\definecolor{DeckMuted}{HTML}{5F6B7A}
\definecolor{DeckSoft}{HTML}{F7F7F5}

\setbeamersize{text margin left=0.55in,text margin right=0.55in}
\setbeamercolor{normal text}{fg=DeckInk,bg=white}
\setbeamercolor{structure}{fg=CalPolyGreen}
\setbeamercolor{title}{fg=white}
\setbeamercolor{frametitle}{fg=CalPolyGreen,bg=white}
\setbeamercolor{block title}{bg=CalPolySage,fg=white}
\setbeamercolor{block body}{bg=CalPolyMint,fg=black}
\setbeamercolor{alerted text}{fg=CalPolyGold}

\setbeamerfont{title}{size=\Huge,series=\bfseries}
\setbeamerfont{frametitle}{size=\Large,series=\bfseries}
\setbeamerfont{block title}{size=\normalsize,series=\bfseries}

\setbeamertemplate{navigation symbols}{}
\setbeamertemplate{itemize item}{\color{CalPolyGold}$\blacktriangleright$}
\setbeamertemplate{itemize subitem}{\color{CalPolyGreen}$\bullet$}
\setbeamertemplate{footline}{
  \hfill{\color{DeckMuted}\scriptsize\insertframenumber/\inserttotalframenumber}
  \hspace{0.18in}\vspace{0.12in}
}
\setbeamertemplate{frametitle}{
  \vspace{0.18in}
  \hspace{0.02in}\insertframetitle\par
  \vspace{0.08in}
  \noindent\makebox[\linewidth][l]{%
    {\color{CalPolyGold}\rule{0.55in}{2pt}}%
    {\color{CalPolyGreen}\rule{\dimexpr\linewidth-0.55in\relax}{2pt}}%
  }
}

\newcommand{\CourseNumber}{Course}
\newcommand{\CourseTitle}{Course Title}
\newcommand{\LectureNumber}{0}
\newcommand{\LectureTitle}{Lecture Title}
\newcommand{\LectureDate}{Date}
\newcommand{\CourseUnit}{Unit}

\newcommand{\coursenumber}[1]{\renewcommand{\CourseNumber}{#1}}
\newcommand{\coursetitle}[1]{\renewcommand{\CourseTitle}{#1}}
\newcommand{\lecturenumber}[1]{\renewcommand{\LectureNumber}{#1}}
\newcommand{\lecturetitle}[1]{\renewcommand{\LectureTitle}{#1}}
\newcommand{\lecturedate}[1]{\renewcommand{\LectureDate}{#1}}
\newcommand{\courseunit}[1]{\renewcommand{\CourseUnit}{#1}}

\author{Kirk Duran}
\institute{California Polytechnic State University, San Luis Obispo}
\date{\LectureDate}

\newcommand{\makedecktitle}{
  \begingroup
  \setbeamercolor{background canvas}{bg=CalPolyGreen}
  \begin{frame}[plain]
    \vfill
    \begin{center}
      {\usebeamerfont{title}\color{white}\LectureTitle\par}
      \vspace{0.18in}
      {\Large\color{white}Lecture \LectureNumber\par}
    \end{center}
    \vfill
    \noindent{\color{white}\rule{\linewidth}{1.2pt}}\par
    \vspace{0.08in}
    \noindent\colorbox{CalPolyGold}{\makebox[0.24\linewidth][c]{\color{white}\Large\CourseNumber}}
    \hspace{0.12in}{\color{white}\Large Kirk Duran}
  \end{frame}
  \endgroup
}

\newenvironment{keyidea}{\begin{block}{Key idea}}{\end{block}}
\newenvironment{activity}{\begin{block}{In class}}{\end{block}}
\newenvironment{cpdefinition}{\begin{block}{Definition}}{\end{block}}
\newenvironment{exampleblockcp}{
  \begingroup
  \setbeamercolor{block title}{bg=CalPolySand,fg=white}
  \setbeamercolor{block body}{bg=CalPolyBeige,fg=black}
  \begin{block}{Example}
}{
  \end{block}
  \endgroup
}

\newcommand{\imageplaceholder}[2][0.3in]{%
  \noindent\makebox[\linewidth][c]{%
    \fcolorbox{CalPolySage}{CalPolyBeige}{%
      \parbox[c][#1][c]{0.86\linewidth}{%
        \centering
        {\color{DeckMuted}\scriptsize Image placeholder: }%
        {\color{DeckInk}\scriptsize\ttfamily\detokenize{#2}}%
      }%
    }%
  }%
}

\newcommand{\sectionframe}[1]{
  \begingroup
  \setbeamercolor{background canvas}{bg=CalPolyGreen}
  \begin{frame}[plain]
    \vfill
    {\color{white}\LARGE\bfseries #1\par}
    \vspace{0.18in}
    {\color{CalPolyGold}\rule{0.22\paperwidth}{3pt}}
    {\color{white}\rule{0.62\paperwidth}{3pt}}
    \vfill
  \end{frame}
  \endgroup
}

\newcommand{\placeholdernote}{
  \begin{block}{Placeholder}
    Replace these bullets with the final lecture material, examples, and in-class activities.
  \end{block}
}

\usepackage[T1]{fontenc}
\usepackage{lmodern}
\pdfstringdefDisableCommands{\def\translate#1{#1}}
\setbeamersize{text margin left=.42in,text margin right=.42in}

% Instructor draft: follows Lecture 31's theme, example/code/trace rhythm,
% fragile verbatim code frames, and shared block environments.
% Suggested 50-minute pacing: invariants 10; Two Sum 12; nondecreasing 7;
% nested-object grouping 8; positive streaks 10; Lab 6 bridge 3.
% Prediction/reveal frames are discussion pauses, not additional topics.
\coursenumber{CS 1001}
\coursetitle{Introduction to Computing}
\lecturenumber{32}
\lecturetitle{Algorithm Design and Loop Invariants}
\lecturedate{November 13, 2026}
\courseunit{12. Algorithms}

\begin{document}
\makedecktitle

\begin{frame}{Today}
\small
\begin{itemize}
 \item Explain why a loop works by describing the meaning of its state.
 \item Design Two Sum using pairs, then remembered information.
 \item Check nondecreasing order using neighboring elements.
 \item Traverse nested objects and group students by major.
 \item Track a current positive streak and the best streak seen so far.
 \item Carry the same reasoning into Lab 6: strings and mixed data.
\end{itemize}
\begin{keyidea}What must the program remember, and what should that state mean after each iteration?\end{keyidea}
\end{frame}

\begin{frame}{From Last Time: A Sorted Region Means Something}
\small
\begin{columns}[T,onlytextwidth]
\begin{column}{.48\textwidth}
\begin{block}{Selection sort}
The finished prefix contains the smallest values in their final sorted positions.
\end{block}
\end{column}
\begin{column}{.48\textwidth}
\begin{block}{Insertion sort}
The processed prefix contains its original values, now in sorted order.
\end{block}
\end{column}
\end{columns}
\begin{itemize}
 \item We used these statements to explain each algorithm's progress.
 \item The statement stays true as the relevant region grows.
 \item Today we will name this reasoning and use it to design new loops.
\end{itemize}
\end{frame}

\begin{frame}{Our Design Routine}
\small
\begin{center}
\textbf{Prompt} $\longrightarrow$ \textbf{Contract} $\longrightarrow$ \textbf{Examples}\\[8pt]
$\longrightarrow$ \textbf{State and invariant} $\longrightarrow$ \textbf{Code} $\longrightarrow$ \textbf{Tests}
\end{center}
\begin{itemize}
 \item Decide what the function receives, returns, and may change.
 \item Work out ordinary and boundary examples before coding.
 \item Give each state variable a precise job.
 \item Explain why one iteration preserves those meanings.
\end{itemize}
\begin{keyidea}Choose the state from the problem, rather than choosing variable names first and hoping they work.\end{keyidea}
\end{frame}

\sectionframe{Part 1: What Makes a Loop Correct?}

\begin{frame}{Definition: Loop Invariant}
\small
\begin{cpdefinition}
A \textbf{loop invariant} is a statement about program state that holds every time execution reaches a specified point in a loop.
\end{cpdefinition}
\begin{block}{Our usual checkpoint}
Immediately before the next iteration: what is already true about the portion processed so far?
\end{block}
\begin{itemize}
 \item State the checkpoint; a claim need not hold halfway through an update.
 \item The variables may change while the relationship remains true.
 \item A useful invariant connects the changing state to the required result.
\end{itemize}
\end{frame}

\begin{frame}[fragile]{A Familiar Invariant: Counting Matches}
\small
\begin{block}{Count values equal to a target}
\footnotesize
\begin{verbatim}
def count_matches(values: list[int], target: int) -> int:
    count = 0
    for value in values:
        if value == target:
            count += 1
    return count
\end{verbatim}
\end{block}
\begin{keyidea}Before each iteration, \texttt{count} equals the number of target matches in the values already processed.\end{keyidea}
\begin{activity}What does the claim say before the first iteration? After the last?\end{activity}
\end{frame}

\begin{frame}{Three Questions About the Invariant}
\small
\begin{enumerate}
 \item \textbf{Initialization:} Why is it true before any iterations?\\
 No values have been processed, so the matching count is \texttt{0}.
 \item \textbf{Maintenance:} Why does one iteration preserve it?\\
 A match adds one; a nonmatch adds nothing. The count now includes the new value.
 \item \textbf{Conclusion at exit:} What follows when processing finishes?\\
 Every value has been processed, so the count describes the entire input.
\end{enumerate}
\begin{keyidea}An invariant explains a whole family of intermediate states, rather than only one example run.\end{keyidea}
\end{frame}

\begin{frame}{Correct State and Eventual Termination}
\small
\begin{columns}[T,onlytextwidth]
\begin{column}{.48\textwidth}
\begin{block}{Does the state stay meaningful?}
Use initialization and maintenance. Combine the invariant with the exit condition to explain the answer.
\end{block}
\end{column}
\begin{column}{.48\textwidth}
\begin{block}{Will execution reach the exit?}
A \texttt{for} loop over our finite, unchanged input eventually exhausts it. A \texttt{while} loop needs progress toward a false condition.
\end{block}
\end{column}
\end{columns}
\begin{itemize}
 \item For binary search, the remaining interval strictly shrinks.
 \item For an early return, explain why the answer is justified at that point.
 \item Passing tests gives evidence; invariant reasoning explains why the updates work generally.
\end{itemize}
\end{frame}

\sectionframe{Part 2: Two Sum}

\begin{frame}{Prompt 1: Find Two Values with a Target Sum}
\small
\begin{activity}
Given a list of integers and a target, return two distinct indexes whose values add to the target. Return \texttt{None} if no such pair exists.
\end{activity}
\begin{center}
\begin{tabular}{lll}
\textbf{Values} & \textbf{Target} & \textbf{One valid result}\\\hline
\texttt{[2, 7, 11, 15]} & \texttt{9} & \texttt{(0, 1)}\\
\texttt{[3, 3]} & \texttt{6} & \texttt{(0, 1)}\\
\texttt{[3]} & \texttt{6} & \texttt{None}
\end{tabular}
\end{center}
\begin{block}{Contract}
Return indexes, not values. Any valid pair is acceptable. Do not modify the input.
\end{block}
\end{frame}

\begin{frame}{Two Sum: Design Before Code}
\small
\begin{activity}
How can we examine every pair without pairing an element with itself?
\end{activity}
\begin{enumerate}
 \item Choose a first index \texttt{i}.
 \item Choose a second index \texttt{j} strictly after \texttt{i}.
 \item Check whether their values add to the target.
 \item Return on a match; report failure only after all pairs are ruled out.
\end{enumerate}
\begin{keyidea}Using \texttt{j > i} guarantees distinct positions and avoids examining the same pair in both orders.\end{keyidea}
\end{frame}

\begin{frame}[fragile]{Two Sum: Examine Pairs Directly}
\small
\begin{block}{First correct solution}
\footnotesize
\begin{verbatim}
def two_sum_pairs(values: list[int], target: int):
    for i in range(len(values)):
        for j in range(i + 1, len(values)):
            if values[i] + values[j] == target:
                return (i, j)
    return None
\end{verbatim}
\end{block}
\begin{itemize}
 \item Before each outer iteration, every pair with a first index less than \texttt{i} has been checked without a match.
 \item The inner loop examines every possible later partner for \texttt{i}.
 \item In the worst case there are $n(n-1)/2$ pairs: quadratic growth.
\end{itemize}
\end{frame}

\begin{frame}{Two Sum: What Information Would Save Work?}
\small
\begin{block}{Suppose the target is 9 and the current value is 7}
The missing value is \texttt{9 - 7}, which is \texttt{2}.
\end{block}
\begin{activity}Can we remember whether \texttt{2} occurred earlier, along with where it occurred?\end{activity}
\begin{itemize}
 \item A dictionary can map a previously seen \textbf{value} to an earlier \textbf{index}.
 \item For each current value, look for \texttt{target - value}.
 \item If it exists, the earlier index and current index form a pair.
 \item Otherwise, remember the current value for later iterations.
\end{itemize}
\begin{keyidea}The dictionary stores exactly the information a future iteration needs.\end{keyidea}
\end{frame}

\begin{frame}[fragile]{Two Sum: Remember Earlier Values}
\small
\begin{block}{Dictionary-based solution}
\footnotesize
\begin{verbatim}
def two_sum(values: list[int], target: int):
    seen = {}
    for i, value in enumerate(values):
        needed = target - value
        if needed in seen:
            return (seen[needed], i)
        seen[value] = i
    return None
\end{verbatim}
\end{block}
\begin{keyidea}Before index \texttt{i} is processed, \texttt{seen} records an earlier index for every distinct value in \texttt{values[:i]}.\end{keyidea}
\end{frame}

\begin{frame}{Trace Two Sum: A Match Can Use Equal Values}
\small
For \texttt{values = [3, 3]} and \texttt{target = 6}:
\begin{center}
\begin{tabular}{cclll}
\textbf{i} & \textbf{value} & \textbf{needed} & \textbf{seen before check} & \textbf{Action}\\\hline
0 & 3 & 3 & \texttt{\{\}} & store \texttt{3: 0}\\
1 & 3 & 3 & \texttt{\{3: 0\}} & return \texttt{(0, 1)}
\end{tabular}
\end{center}
\begin{activity}What happens for \texttt{[3]} with the same target? What if we store the current value before checking?\end{activity}
\begin{keyidea}Check first, then store. Earlier indexes are distinct from the current index; equal values are still allowed.\end{keyidea}
\end{frame}

\begin{frame}{Why the Dictionary Solution Works}
\small
\begin{enumerate}
 \item \textbf{Initialize:} the processed prefix is empty and \texttt{seen} is empty.
 \item \textbf{Maintain:} after an unsuccessful check, record the current value's index. Every distinct processed value still has a representative.
 \item \textbf{Successful exit:} the saved index is earlier, and its value plus the current value equals the target.
 \item \textbf{Exhaustion:} if a valid pair existed, its later member would have found the earlier member's value. No pair was missed.
\end{enumerate}
\begin{block}{Work and storage}
One traversal; under the usual average constant-time dictionary-operation model, expected $O(n)$ time and up to $O(n)$ extra storage. The pair solution uses $O(1)$ extra storage.
\end{block}
\end{frame}

\begin{frame}{Two Sum: Predict and Test}
\small
\begin{activity}Predict the dictionary solution's return value for each case.\end{activity}
\begin{center}
\begin{tabular}{lll}
\textbf{Values} & \textbf{Target} & \textbf{What this checks}\\\hline
\texttt{[]} & 4 & no iterations\\
\texttt{[4]} & 8 & cannot reuse one position\\
\texttt{[4, 4]} & 8 & repeated values, distinct positions\\
\texttt{[-3, 8, 2]} & 5 & negative values\\
\texttt{[1, 2, 4]} & 8 & no pair\\
\texttt{[1, 4, 2, 3]} & 5 & multiple possible answers
\end{tabular}
\end{center}
\begin{keyidea}A general correctness test checks the returned indexes and their sum; different correct algorithms may choose different pairs.\end{keyidea}
% Reveal in row order: None; None; (0, 1); (0, 1); None; (0, 1).
\end{frame}

\sectionframe{Part 3: Neighboring Values}

\begin{frame}{Prompt 2: Is the List Nondecreasing?}
\small
\begin{activity}
Return \texttt{True} exactly when each value is at least as large as its predecessor. Leave the list unchanged.
\end{activity}
\begin{center}
\begin{tabular}{ll}
\textbf{Input} & \textbf{Result}\\\hline
\texttt{[2, 2, 5, 9]} & \texttt{True}\\
\texttt{[2, 5, 3, 9]} & \texttt{False}\\
\texttt{[]} & \texttt{True}\\
\texttt{[7]} & \texttt{True}
\end{tabular}
\end{center}
\begin{block}{Contract detail}Equal neighbors are allowed. Empty and singleton lists have no violating adjacent pair.\end{block}
\end{frame}

\begin{frame}{Nondecreasing: Choose the Comparisons}
\small
\begin{center}
\texttt{[2, 2, 5, 9]}\\[8pt]
\texttt{2 <= 2}\qquad\texttt{2 <= 5}\qquad\texttt{5 <= 9}
\end{center}
\begin{activity}
Which indexes identify these pairs? Why should the current index start at \texttt{1}?
\end{activity}
\begin{itemize}
 \item Compare \texttt{values[i - 1]} with \texttt{values[i]}.
 \item One decrease is enough to disprove the whole claim.
 \item One correctly ordered pair is not enough to prove the whole claim.
\end{itemize}
\begin{keyidea}Before checking index \texttt{i}, every adjacent pair ending before \texttt{i} has passed.\end{keyidea}
\end{frame}

\begin{frame}[fragile]{Nondecreasing: Implementation}
\small
\begin{block}{Look for a counterexample}
\footnotesize
\begin{verbatim}
def is_nondecreasing(values: list[int]) -> bool:
    for i in range(1, len(values)):
        if values[i] < values[i - 1]:
            return False
    return True
\end{verbatim}
\end{block}
\begin{itemize}
 \item Before the first check, no adjacent pairs need to have passed yet.
 \item A successful comparison extends the checked prefix by one value.
 \item Reaching the final return means every adjacent pair passed.
 \item For fewer than two values, the loop is empty and returns \texttt{True}.
\end{itemize}
\end{frame}

\begin{frame}{Nondecreasing: Trace and Repair}
\small
\begin{center}
\begin{tabular}{clll}
\textbf{i} & \textbf{Previous} & \textbf{Current} & \textbf{Action for [2, 5, 3, 9]}\\\hline
1 & 2 & 5 & no decrease; continue\\
2 & 5 & 3 & decrease; return \texttt{False}\\
3 & --- & --- & never reached
\end{tabular}
\end{center}
\begin{activity}
What goes wrong if \texttt{return True} is inside the loop, immediately after the \texttt{if} statement?
\end{activity}
\begin{keyidea}A successful first comparison does not establish the property for the unchecked suffix.\end{keyidea}
\begin{block}{Work}At most $n-1$ comparisons for $n\geq1$: $O(n)$ worst-case time and $O(1)$ extra storage. We do not need to sort the input to check it.\end{block}
\end{frame}

\sectionframe{Part 4: Algorithms over Nested Objects}

\begin{frame}{Our Shared Object Model}
\small
\begin{center}
\begin{tabular}{cl}
\textbf{Object} & \textbf{Attributes}\\\hline
\texttt{Course} & \texttt{code}, \texttt{students: list[Student]}\\[6pt]
$\downarrow$ & the roster contains references to student objects\\[6pt]
\texttt{Student} & \texttt{name}, \texttt{major}, \texttt{scores: list[int]}\\[6pt]
$\downarrow$ & scores are in chronological order
\end{tabular}
\end{center}
\begin{block}{Contracts for today's examples}
Each student appears once in the roster, names are unique within a course, and every score list is finite. Algorithms read the course and students without changing them.
\end{block}
\begin{keyidea}Follow the data one level at a time: course, student, then score.\end{keyidea}
\end{frame}

\begin{frame}[fragile]{Define the Objects We Will Use}
\small
\begin{columns}[T,onlytextwidth]
\begin{column}{.48\textwidth}
\begin{block}{Student state}
\footnotesize
\begin{verbatim}
class Student:
    def __init__(self, name: str,
                 major: str,
                 scores: list[int]):
        self.name = name
        self.major = major
        self.scores = scores
\end{verbatim}
\end{block}
\end{column}
\begin{column}{.48\textwidth}
\begin{block}{A course contains students}
\footnotesize
\begin{verbatim}
class Course:
    def __init__(self, code: str,
                 students: list[Student]):
        self.code = code
        self.students = students
\end{verbatim}
\end{block}
\end{column}
\end{columns}
\begin{keyidea}These assignments store references. The constructors do not automatically copy the roster or score lists.\end{keyidea}
\end{frame}

\begin{frame}[fragile]{Construct One Course and Three Students}
\small
\begin{block}{Shared data for both object exercises}
\footnotesize
\begin{verbatim}
ada = Student("Ada", "CS", [2, 3, 0, 4, 5, 6, 0])
grace = Student("Grace", "Math", [0, 5, 6])
linus = Student("Linus", "CS", [])
course = Course("CS 1001", [ada, grace, linus])
\end{verbatim}
\end{block}
\begin{activity}Predict each expression before running it.\end{activity}
\begin{center}
\begin{tabular}{ll}
\texttt{course.students[0] is ada} & \texttt{?}\\
\texttt{course.students[1].major} & \texttt{?}\\
\texttt{course.students[0].scores[3]} & \texttt{?}
\end{tabular}
\end{center}
% Reveal: True; "Math"; 4. Access the roster, then the student, then its scores.
\end{frame}

\begin{frame}{Prompt 3: Group a Course's Students by Major}
\small
\begin{activity}
Given a \texttt{Course}, return a new dictionary mapping each major to a new list of that major's students. Preserve roster order within each group.
\end{activity}
\begin{block}{For our example course}
\texttt{\{"CS": [ada, linus], "Math": [grace]\}}
\end{block}
\begin{itemize}
 \item The values in each group are the original student objects.
 \item The dictionary and its group lists are new containers.
 \item An empty roster produces an empty dictionary.
\end{itemize}
\begin{activity}What does a group store that a frequency counter does not?\end{activity}
\end{frame}

\begin{frame}[fragile]{Grouping: Plan and Implementation}
\small
\begin{block}{One group for each observed major}
\footnotesize
\begin{verbatim}
def group_by_major(course: Course) -> dict[str, list[Student]]:
    groups = {}
    for student in course.students:
        major = student.major
        if major not in groups:
            groups[major] = []
        groups[major].append(student)
    return groups
\end{verbatim}
\end{block}
\begin{keyidea}After each iteration, every processed student appears exactly once in the correct group, in roster order.\end{keyidea}
Start empty; create each bucket once; append each student once.
\end{frame}

\begin{frame}{Grouping: Trace the Changing Containers}
\small
\begin{center}
\begin{tabular}{lll}
\textbf{Student} & \textbf{Major} & \textbf{groups after processing}\\\hline
--- & --- & \texttt{\{\}}\\
Ada & CS & \texttt{\{"CS": [ada]\}}\\
Grace & Math & \texttt{\{"CS": [ada], "Math": [grace]\}}\\
Linus & CS & \texttt{\{"CS": [ada, linus], "Math": [grace]\}}
\end{tabular}
\end{center}
\begin{activity}
Where is each new list created? Why must \texttt{groups[major] = []} run only when the major is first encountered?
\end{activity}
\begin{keyidea}Reinitializing an existing group's list would discard earlier students from that group.\end{keyidea}
\end{frame}

\begin{frame}[fragile]{New Containers Can Still Share Their Elements}
\small
\begin{block}{Predict each effect}
\footnotesize
\begin{verbatim}
groups = group_by_major(course)
print(groups["CS"][0] is course.students[0])
groups["CS"].pop()
print(len(course.students))
groups["CS"][0].name = "Ada Lovelace"
print(course.students[0].name)
\end{verbatim}
\end{block}
\begin{activity}Which statements change a container? Which change a shared student?\end{activity}
% Reveal: True; 3; Ada Lovelace. The grouping algorithm itself does not mutate inputs.
% This caller experiment mutates the group list and then a shared student.
% Restore ada.name = "Ada" before continuing with the shared course example.
\end{frame}

\begin{frame}{Reference Exercise: Reveal}
\small
\begin{center}
\begin{tabular}{ll}
\textbf{Output} & \textbf{Reason}\\\hline
\texttt{True} & both containers refer to the original Ada object\\[6pt]
\texttt{3} & removing from a group does not remove from the roster\\[6pt]
\texttt{Ada Lovelace} & attribute assignment changes the shared student
\end{tabular}
\end{center}
\begin{keyidea}A fresh outer container does not imply fresh copies of the objects inside it.\end{keyidea}
\begin{block}{Reset for the remaining examples}Restore \texttt{ada.name = "Ada"}. The grouping function itself left the original objects unchanged; the caller performed these later mutations.\end{block}
\end{frame}

\sectionframe{Part 5: Current Streak and Best Streak}

\begin{frame}{Prompt 4: A Student's Longest Positive Streak}
\small
\begin{activity}
Add a method that returns the longest consecutive run of scores greater than zero. Leave the student's scores unchanged.
\end{activity}
\begin{center}
\begin{tabular}{ll}
\textbf{Scores} & \textbf{Longest positive streak}\\\hline
\texttt{[2, 3, 0, 4, 5, 6, 0]} & \texttt{3}\\
\texttt{[0, 5, 6]} & \texttt{2}\\
\texttt{[]} & \texttt{0}\\
\texttt{[-2, 0, -1]} & \texttt{0}
\end{tabular}
\end{center}
\begin{activity}Why does counting all positive scores answer a different question?\end{activity}
\end{frame}

\begin{frame}{Streaks: Give Two Variables Different Jobs}
\small
\begin{columns}[T,onlytextwidth]
\begin{column}{.48\textwidth}
\begin{block}{\texttt{current}}
The number of consecutive positive scores at the \textbf{end} of the processed prefix.
\end{block}
\end{column}
\begin{column}{.48\textwidth}
\begin{block}{\texttt{best}}
The longest positive streak \textbf{anywhere} in the processed prefix.
\end{block}
\end{column}
\end{columns}
\begin{itemize}
 \item A positive score extends \texttt{current}; any other score resets it.
 \item Update \texttt{best} if the current streak exceeds it.
 \item A reset of \texttt{current} must not erase \texttt{best}.
\end{itemize}
\begin{keyidea}At the iteration boundary, both meanings must hold. Neither variable alone represents both facts.\end{keyidea}
\end{frame}

\begin{frame}[fragile]{Streaks: Add a Method to Student}
\small
\begin{block}{Place this method inside the existing Student class}
\footnotesize
\begin{verbatim}
    def longest_positive_streak(self) -> int:
        current = 0
        best = 0
        for score in self.scores:
            if score > 0:
                current += 1
            else:
                current = 0
            if current > best:
                best = current
        return best
\end{verbatim}
\end{block}
\begin{keyidea}Reads \texttt{self.scores}; returns a number without mutating the student.\end{keyidea}
\end{frame}

\begin{frame}{Streaks: Trace Ada's Scores}
\small
\begin{center}
\begin{tabular}{clcc}
\textbf{Score} & \textbf{Action} & \textbf{current after update} & \textbf{best after update}\\\hline
--- & initialize & 0 & 0\\
2 & extend & 1 & 1\\
3 & extend & 2 & 2\\
0 & reset & 0 & 2\\
4 & extend & 1 & 2\\
5 & extend & 2 & 2\\
6 & extend & 3 & 3\\
0 & reset & 0 & 3
\end{tabular}
\end{center}
\begin{activity}What would change if the scores ended at \texttt{6}? Why does returning \texttt{current} fail for the full example?\end{activity}
\end{frame}

\begin{frame}{Streaks: Explain Correctness}
\small
\begin{enumerate}
 \item \textbf{Initialize:} an empty prefix has an ending streak of length zero and a best streak of length zero.
 \item \textbf{Maintain current:} a positive score extends the ending streak; a nonpositive score makes its length zero.
 \item \textbf{Maintain best:} the only new candidate for a longer streak ends at the new score. Compare its length with the previous best.
 \item \textbf{At exit:} the processed prefix is the full score list, so \texttt{best} is the required result.
\end{enumerate}
\begin{block}{Termination and work}
The method traverses one finite score list: $O(m)$ time for $m$ scores and $O(1)$ extra storage.
\end{block}
\end{frame}

\begin{frame}[fragile]{Now Report the Streak for Every Student}
\small
\begin{activity}Return a dictionary mapping each student's unique name to their longest positive streak.\end{activity}
\begin{block}{The outer traversal follows the object structure}
\footnotesize
\begin{verbatim}
def streak_report(course: Course) -> dict[str, int]:
    report = {}
    for student in course.students:
        report[student.name] = student.longest_positive_streak()
    return report
\end{verbatim}
\end{block}
For our example: \texttt{\{"Ada": 3, "Grace": 2, "Linus": 0\}}.
\begin{keyidea}The outer loop visits students; each method call traverses that student's scores and starts with fresh local accumulators.\end{keyidea}
\end{frame}

\begin{frame}{Nested Traversal: State, Work, and Boundary Cases}
\small
\begin{itemize}
 \item \textbf{Outer invariant:} the report contains the correct result for every student already processed.
 \item \textbf{Inner invariant:} \texttt{current} and \texttt{best} describe the processed scores of one student.
 \item Each method call initializes its own \texttt{current} and \texttt{best}; one student's streak does not carry into another's.
 \item An empty course gives \texttt{\{\}}; a student with no scores gets \texttt{0}.
 \item With $s$ students and $q$ total scores, the work is $O(s+q)$; students can have different numbers of scores.
\end{itemize}
\begin{activity}Why would a report keyed by name lose information if two students had the same name?\end{activity}
% A later assignment could use student IDs; today unique names are a precondition.
\end{frame}

\begin{frame}{Streaks: Predict Before Running}
\small
\begin{center}
\begin{tabular}{ll}
\textbf{Scores} & \textbf{What this checks}\\\hline
\texttt{[]} & no iterations\\
\texttt{[5]} & one positive score\\
\texttt{[0, -2]} & no positive score\\
\texttt{[1, 2, 3]} & all positive; best streak at the end\\
\texttt{[1, 2, 0, 3]} & best streak before the end\\
\texttt{[1, 0, 2, 0, 3]} & separated positives\\
\texttt{[1, 2, 0, 3, 4]} & tied longest streaks
\end{tabular}
\end{center}
\begin{activity}Give each expected result. Then choose two students whose scores would expose a failure to reset state between students.\end{activity}
% Reveal: 0; 1; 0; 3; 2; 1; 2. Try one student with [1, 2, 3]
% followed by a student with [1], whose result must be 1, not 4 or 3.
\end{frame}

\sectionframe{Part 6: Carry the Reasoning into Lab 6}

\begin{frame}{Lab 6: Different Data, Familiar State Meanings}
\small
\begin{center}
\begin{tabular}{@{}lp{.65\linewidth}@{}}
\textbf{Task} & \textbf{What should be true of the processed prefix?}\\\hline
\texttt{count\_character} & The count equals the matching characters seen so far.\\[5pt]
\shortstack[l]{\texttt{swap\_case},\\\texttt{str\_translate}} & The new string is the required transformation of the processed characters.\\[5pt]
\shortstack[l]{\texttt{remove\_character},\\\texttt{words\_of\_length}} & The result contains exactly the selected characters or words, in encounter order.\\[5pt]
\texttt{groups\_of\_3} & The groups built so far, plus any unfinished group, account for every processed value exactly once, in order.
\end{tabular}
\end{center}
\begin{keyidea}Strings are immutable: construct a new result. For nested lists, give each output group its own container.\end{keyidea}
\end{frame}

\begin{frame}{Lab 6: Turn the Contract into Tests}
\small
\begin{itemize}
 \item Use \texttt{unittest.TestCase} and \texttt{self.assertEqual(...)} for your added cases.
 \item Test empty inputs, one element, first/last matches, no matches, and repeated matches.
 \item Check that transformations and filters preserve the required order.
 \item For \texttt{groups\_of\_3}, test lengths divisible by three and final groups of one or two.
 \item Check that the input list remains unchanged and distinct output groups do not share one mutable list.
 \item Follow the lab's restrictions on shortcut string methods.
\end{itemize}
\begin{activity}For an eight-element input to \texttt{groups\_of\_3}, what must still be handled after the last full group?\end{activity}
% Expected: the last two elements must appear once in a final two-element group.
\end{frame}

\begin{frame}{Exit Ticket: Explain the State}
\small
\begin{enumerate}
 \item Why must Two Sum check the dictionary before storing the current value?
 \item Why can nondecreasing return \texttt{False} early but not \texttt{True} after just one successful comparison?
 \item Why do \texttt{current} and \texttt{best} need different meanings?
 \item If two output groups share one list object, what happens when one group is appended to?
\end{enumerate}
\begin{keyidea}An explanation should connect the update to the meaning of the state, rather than only repeat the Python statement.\end{keyidea}
\end{frame}

\begin{frame}{Takeaways}
\small
\begin{enumerate}
 \item A loop invariant states what is true at a specified iteration boundary.
 \item Explain initialization, maintenance, the conclusion at exit, and termination.
 \item A dictionary can remember information that avoids repeated searching.
 \item Adjacent comparisons answer questions about neighboring values.
 \item Nested objects require following references and restarting local state at the right level.
 \item Fresh containers can hold shared objects; distinguish container changes from element changes.
 \item Give each accumulator a precise job, then test the boundaries of that meaning.
\end{enumerate}
\begin{block}{Next}Input, output, and string processing. In Lab 6, apply today's reasoning to strings and nested lists.\end{block}
\end{frame}

\end{document}
